Una empresa fabrica el mismo modelo de coche en cuatro fábricas diferentes. Una vez fabricados los coches se transportan a una campa que existe en cada uno de lo los cinco continentes en los que la empresa vende su vehículo.

Las capacidades (en cientos de miles de coches) de cada fábrica son 10, 20, 15 y 20. Las demandas de cada continente son respectivamente 15, 10, 15, 15 y 10.

La siguiente tabla indica los costes de transporte de un coche desde una fábrica (fila) a un continente (columna).

\begin{table}[h]
\begin{center}
\begin{tabular}{|c|c|c|c|c|}
\hline
4 & 2 & 3 & 5 & 1 \\ \hline
5 & 2 & 7 & 3 & 8 \\ \hline
6 & 4 & 4 & 3 & 7 \\ \hline
1 & 2 & 5 & 3 & 4 \\ \hline
\end{tabular}
\end{center}
\end{table}

\begin{enumerate}
    \item \textbf{Aplica} el método de Vogel para encontrar una solución inicial de este problema de transporte.
   
    
    \item ¿Qué \textbf{características} tiene la solución que se ha obtenido aplicando el método de Vogel?¿Por qué?

     
    El problema está equilibrado.

    Una posible solución al aplicar el método de Vogel es ésta (hay más posibles dados los empates que aparecen):

\begin{table}[h!]
\begin{center}
\begin{tabular}{|l|l|l|l|l|}
\hline
   &    &    &    & 10 \\ \hline
   & 10 &    & 10 &    \\ \hline
   &    & 15 &    &    \\ \hline
15 &    &    & 5  &    \\ \hline
\end{tabular}
\end{center}
\end{table}
    Esta solución es:
    \begin{itemize}
        \item básica, el número de variables no nulas es menor de 4+5-1=8.
        \item factible, cumple las 9 restricciones.
        \item degenerada, sólo tiene 6 variables básicas no nulas.
        \item óptima, no existen costes reducidos positivos.
    \end{itemize}
\end{enumerate}

